\section{Ringkasan Konteks Kurikulum OBE}

Pendekatan \textit{Outcome-Based Education} (OBE) menggeser fokus pendidikan dari "apa yang diajarkan" menjadi "apa yang mampu dilakukan oleh mahasiswa" setelah proses pembelajaran selesai. Dalam sistem ini, keberhasilan pendidikan tidak diukur dari berapa banyak jam materi yang tersampaikan, melainkan dari bukti nyata capaian kompetensi yang dapat ditunjukkan oleh mahasiswa pada akhir periode belajar \cite{obe_wikipedia}.

Oleh karena itu, \textit{learning outcomes} dalam buku ini disusun menggunakan kata kerja operasional yang spesifik dan terukur agar dapat diamati serta diases secara objektif \cite{adelaide_learning_outcomes,stanford_learning_outcomes}. Dalam konteks mata kuliah Kriptografi, hirarki kurikulum yang diterapkan adalah sebagai berikut:

\begin{enumerate}
    \item \textbf{CPL (Capaian Pembelajaran Lulusan)}: Merupakan profil kompetensi akhir yang harus dimiliki oleh setiap lulusan Program Studi Teknik Informatika. Ini adalah standar kualitas lulusan yang mencakup aspek sikap, pengetahuan, dan keterampilan.
    \item \textbf{CPMK (Capaian Pembelajaran Mata Kuliah)}: Merupakan turunan dari CPL yang menjadi target utama mata kuliah Kriptografi. Contohnya adalah kemampuan \textit{mengenal} algoritma kriptografi (CPMK-1) dan \textit{memahami} urgensi keamanan informasi (CPMK-2).
    \item \textbf{Sub-CPMK}: Merupakan indikator pencapaian yang lebih detail dan spesifik. Sub-CPMK inilah yang menjadi acuan dalam penyusunan setiap bab, aktivitas praktikum, soal latihan, dan instrumen asesmen di dalam buku ini.
\end{enumerate}

Implementasi OBE dalam buku ini menuntut adanya keterhubungan yang kuat antara aktivitas pembelajaran dan indikator capaian. Buku ini menekankan penggunaan \textbf{asesmen formatif} (latihan dan aktivitas di tengah bab) serta \textbf{refleksi diri} (checklist kompetensi), sehingga mahasiswa dapat memantau kemajuan belajarnya secara mandiri, mengenali kesenjangan kompetensi yang masih ada, dan mengambil langkah perbaikan secara proaktif sebelum menghadapi asesmen sumatif (ujian).
